\SourceLeaf{60}
$(1+n\cos\varphi)^{\nu-1}$, can be derived. In this way we can pass from $(1+n\cos\varphi)^{-1}$ to higher negative powers, avoiding an infinite series for each coefficient.

$\alpha$) To investigate the law that holds when $\nu=-1$ for the initial coefficient $A$, let us take the general expression for its $h^{\text{ti}}$ term and set $\nu=-1$ in it. Since
\[
{}^{-1}\mathfrak{B}^{(2h)} = 1, \qquad \text{that term} \qquad = \frac{1}{2^{2h}}\,{}^{2h}\mathfrak{B}^{(h)}\,n^{2h}.
\]
Now
\[
\frac{1}{2^{2h}}\,{}^{2h}\mathfrak{B}^{(h)} = \frac{1}{2^{2h}}\,\frac{2h.(2h-1).\,.\,.\,.(h+1)}{1.2.\,.\,.\,. h}
\]
\[
= \frac{1}{2^{2h}} . \frac{2h.(2h-1).\,.\,.\,. 1}{1.2.\,.\,.\,.\,. h} . \frac{1}{h.(h-1).\,.\,. 1}
\]
\[
= \frac{1}{2^{2h}}\,\frac{(2h-1)(2h-3).\,.\,.\,. 1}{1.2.\,.\,. h} . \frac{2h.(2h-2).\,.\,. 2}{h.(h-1).\,.\,.\,.\,. 1}
\]
\[
= \frac{1}{2^{h}} . \frac{(2h-1)(2h-3).\,.\,. 1}{1.2.\,.\,.\,. h},
\]
and through these transformations that $h^{\text{tus}}$ term becomes
\[
= \frac{1}{2^{h}}\,\frac{(2h-1).(2h-3).\,.\,. 1}{1.2.\,.\,.\,. h} .  n^{2h}.
\]
Since this expression is clearly what the series expansion of $(1-nn)^{-\frac{1}{2}}$ gives as its $h^{\text{tus}}$ term, it follows that $A$ itself, when $\nu=-1$, will be
\[
= (1-nn)^{-\frac{1}{2}}
\]
\par\noindent\hspace*{\parindent}$\beta$) The first coefficient after the initial one, $A^{(1)}$, is determined by the relation established earlier (7, $\delta$)
\[
A^{(1)} = \frac{2.(\nu+2)}{n}\int An\,\mathrm{d}n - 2n.A,
\]
which here, because $\nu=-1$, becomes
\[
A^{(1)} = \frac{2}{n}\int An\,\mathrm{d}n - 2nA = \frac{2}{n}\int(1-nn)^{-\frac{1}{2}}n\,\mathrm{d}n - 2n(1-nn)^{-\frac{1}{2}}.
\]
Carrying out the integration (in such a way that the integral vanishes with $n$) gives
\[
A^{(1)} = -\frac{2}{n}(1-nn)^{\frac{1}{2}} + \frac{2}{n} - 2n(1-nn)^{-\frac{1}{2}},
\]
or more briefly
\[
A^{(1)} = \frac{2}{n}\bigl(1-(1-nn)^{-\frac{1}{2}}\bigr).
\]

